% FILE: 01_research_question.tex
% Project: audits
% Role: completeness-audit-and-compliance-verification-corpus

\section{Research Question}
\label{sec:audits-rq}

\subsection{Problem Statement}
\label{sec:audits-rq-problem}

The central problem addressed by the audits corpus is the \emph{verification gap}
between declared manufacturing pipelines and the processes that actually execute.
In any sufficiently complex software system, it is possible~--- and in practice
common~--- for the declared architecture to diverge from the runtime behavior.
Stage outputs may be emitted outside lawful object lifecycles. Proof gates may
pass despite non-conforming execution. Receipts may carry type names that suggest
structured evidence while actually collapsing that evidence to a casual carrier.

For a research foundry that claims to manufacture process-intelligence artifacts
under Van der Aalst process law, this verification gap is not a quality concern
but an existential one. If the foundry's own execution cannot be audited against
its own declared process model, then every downstream claim the foundry makes
about enterprise process conformance is unsupported.

The problem is made concrete by the \textsc{dto\_001} violation discovered in
Audit~1: two method names in \texttt{wasm4pm-compat}~--- \texttt{to\_json\_string()}
at \texttt{compat/src/manufacturing/mod.rs:735} and \texttt{receipt\_json()} at
\texttt{compat/src/manufacturing/traits.rs:34}~--- use the \texttt{\_json} suffix
pattern, indicating that structured \texttt{Evidence<T>} types are being collapsed
to a casual JSON carrier. This is a type-law violation that blocks graduation. It
is exactly the kind of defect that would be invisible to a conventional test suite,
which checks output correctness but not law conformance.

\subsection{Old-Regime Assumption Challenged}
\label{sec:audits-rq-old-regime}

The old-regime assumption challenged by this corpus is: \emph{a passing test suite
is sufficient evidence that a system behaves as declared.}

This assumption underlies most software quality programs. If all unit tests pass,
if the build completes, if the CI pipeline is green, the system is considered
conformant. No further verification of the process that produced those results is
required.

The Van der Aalst Constitution, as encoded in
\texttt{\textasciitilde/.claude/rules/process-mining-chicago-tdd.md}, refuses this
assumption directly: \emph{``Don't trust code paths, state machines, or API
responses. Trust only event evidence that can be mined into a conforming object-
centric process.''} The audits corpus operationalizes this refusal. A system that
compiles and passes tests but contains \texttt{\_json} method names has failed the
process law, regardless of what the test outputs report.

\subsection{Hypothesis}
\label{sec:audits-rq-hypothesis}

The hypothesis of the audits corpus is that deterministic, pattern-level audits of
source artifacts~--- applied at the type-law surface rather than the behavioral
surface~--- can detect process-law violations that passing test suites cannot, and
that these audits can be expressed as reproducible shell programs whose output
constitutes a cryptographically anchored receipt chain.

Concretely: five shell audits targeting DTO flattening, tool smuggling, feature
isolation, projection receipt, and graduation boundary constitute a necessary and
sufficient verification program for the \texttt{wasm4pm-compat} projection layer.
A system that passes all five audits has been confirmed, at the type-law surface,
to be free of the violation classes that would prevent graduation to the
\texttt{wasm4pm} execution engine.

The hypothesis further holds that the same methodology~--- event-log derivation
from source artifacts, model-vs-log conformance checking, temporal ordering of
audit receipts~--- applies recursively to the research foundry itself, making the
foundry self-auditing.

\subsection{Success Criteria}
\label{sec:audits-rq-success}

Success requires all of the following to hold simultaneously:

\begin{enumerate}
  \item All five deterministic shell audits return \textsc{pass} (currently 4/5;
    the \textsc{dto\_001} violation at \texttt{mod.rs:735} and
    \texttt{traits.rs:34} must be remediated).
  \item The \texttt{PROCESS\_INTELLIGENCE\_COMPAT\_CONFORMANCE\_001} checkpoint
    file is committed and signed.
  \item The compile-fail fixture corpus maintains 97.7\% or higher valid
    law-proving error code coverage with zero \texttt{E0425} defects.
  \item The \texttt{ALIVE\_001} gate remains certified: 567 or more commits,
    all 12 corpus directories at or above minimum document count thresholds.
  \item All five residual lifecycle gaps (RESIDUAL-DESIGN-001, RESIDUAL-SIM-001,
    RESIDUAL-SIM-002, RESIDUAL-MON-001, RESIDUAL-MON-002) are closed with
    documented wasm4pm execution boundary mappings.
  \item Board claim traceability is verified across all nine claim types in the
    four strategic domains (EBITDA Optimization, Working Capital, GRC
    Defensibility, Integration Velocity).
\end{enumerate}
